\subsection{Indirect Support Prioritizes Drug Targets With Higher Clinical Success}

\begin{figure}[t!]
    \centering
    {{canva:5}}
    \caption{Drug target validation. (a) Relative success (RS) as a function of the fraction of each resource prioritized, comparing PIGEAN indirect and combined support to Open Targets Genetics (OTG) L2G, Open Targets Platform (OTP) genetic support, and OTP indirect support. (b) Temporal holdout: RS for post-2018 drug targets using pre-2018 geneset libraries.}
    \label{fig:validation-pharma}
\end{figure}

We evaluated whether genes with high indirect support are enriched for successful drug targets. We matched PIGEAN gene--trait associations to drug target--indication (T-I) pairs from the Pharmaprojects clinical trial database across {{NUM_TOTAL_TRAITS_MOUSE_MSIGDB}} phenotypes using MeSH-based semantic similarity (similarity threshold $\geq 0.8$; Section~\ref{sec:methods-rs}). For each T-I pair, we computed the cumulative relative success (RS)---the product of the ratios of phase-advancement probabilities for prioritized versus unprioritized targets across all clinical stages (Preclinical through Launched; Equations~\ref{eq:rs-phase}--\ref{eq:rs-cum})---following the framework of Minikel et al.\ \cite{minikel_refining_2024}. Because PIGEAN, unlike binary classifiers, assigns continuous scores to gene--trait associations, we evaluated RS across the full range of score thresholds using the continuous RS curve described in Section~\ref{sec:methods-rs}. At the RS level achieved by Open Targets Genetics (OTG) with a minimum locus-to-gene (L2G) score of 0.1, as curated by Minikel et al.\ \cite{minikel_refining_2024} (RS = {{MINIKEL_OTG_RS:.2f}}; 95\% CI: {{MINIKEL_OTG_RS_CI_LOWER:.2f}} to {{MINIKEL_OTG_RS_CI_UPPER:.2f}}), PIGEAN's indirect support prioritized {{NUM_PIGEAN_PRIORITIZED_TARGETS_MOUSE_MSIGDB_AT_RS_eq_MINIKEL_OTG_INDIRECT}} targets compared with only {{NUM_MINIKEL_OTG_PRIORITIZED_TARGETS}} from OTG. Combined genetic support (direct plus indirect) at the same RS level prioritized {{NUM_PIGEAN_PRIORITIZED_TARGETS_MOUSE_MSIGDB_AT_RS_eq_MINIKEL_OTG_COMBINED}} targets, with {{NUM_UNIQUE_PIGEAN_INDIRECT_TARGETS_COMPARED_WITH_COMBINED}} targets uniquely identified by indirect support alone, indicating that indirect and combined support surface partially distinct sets of drug targets.

We benchmarked PIGEAN against OTG L2G scores, the Open Targets Platform (OTP) harmonic mean genetic support score, and OTP's literature- and model-based indirect support metric across the full range of thresholds (Figure~\ref{fig:validation-pharma}a). Across this range, PIGEAN's indirect and combined support consistently yielded higher RS than OTP's indirect and genetic support metrics while prioritizing a comparable number of targets. L2G alone performed almost equivalently to PIGEAN's indirect support, but combined PIGEAN support substantially outperformed all approaches over the first 1,000 prioritized targets (Extended Data Figure~\ref{fig:extended-1}). Notably, RS for L2G remained relatively stable across L2G scores, such that even scores close to 1 did not yield markedly increased RS, consistent with the findings of Minikel et al.\ By contrast, both indirect and combined support from PIGEAN achieved substantially higher RS at higher scores, indicating greater selectivity at stringent thresholds.

% Temporal holdout
To control for potential information leakage from gene sets that encode drug target knowledge, we performed a temporal holdout analysis (Section~\ref{sec:methods-rs}): PIGEAN was re-run using only gene set libraries curated from pre-2018 data, and evaluation was restricted to T-I pairs for drugs launched after 2018 (Figure~\ref{fig:validation-pharma}b). Even under this temporal holdout, indirect support achieved a maximum RS of {{MAX_RS_INDIRECT_TEMPORAL_HOLDOUT}} (95\% CI: {{MAX_RS_INDIRECT_TEMPORAL_HOLDOUT_CI_LOWER}} to {{MAX_RS_INDIRECT_TEMPORAL_HOLDOUT_CI_UPPER}}; $N$ = {{NUM_DRUG_TARGETS_SUPPORTED_INDIRECT_TEMPORAL_HOLDOUT}}) and combined support achieved {{MAX_RS_COMBINED_TEMPORAL_HOLDOUT}} (95\% CI: {{MAX_RS_COMBINED_TEMPORAL_HOLDOUT_CI_LOWER}} to {{MAX_RS_COMBINED_TEMPORAL_HOLDOUT_CI_UPPER}}; $N$ = {{NUM_DRUG_TARGETS_SUPPORTED_COMBINED_TEMPORAL_HOLDOUT}}). At 2\% of the resource, RS stabilized at {{RS_INDIRECT_TEMPORAL_HOLDOUT_STABLE_AT_ALPHA_002}} for indirect support and {{RS_COMBINED_TEMPORAL_HOLDOUT_STABLE_AT_ALPHA_002}} for combined support, confirming that the enrichment for successful drug targets is not driven by circularity with existing drug target annotations.
